% Part: computability
% Chapter: recursive-functions
% Section: examples

\documentclass[../../../include/open-logic-section]{subfiles}

\begin{document}

\olfileid[hi]{cmp}{rec}{exa}
\olsection{आदिम पुनरावर्ती फलनों के उदाहरण}


हमारे पास आदिम पुनरावर्ती फलनों के कुछ उदाहरण पहले से हैं: जोड़ और गुणा
फलन~$\Add$ तथा $\Mult$। तत्समक फलन $\fn{id}(x) = x$ आदिम पुनरावर्ती
है, क्योंकि वह केवल~$\Proj{1}{0}$ है। अचर फलन $\fn{const}_n(x) = n$
  आदिम पुनरावर्ती है, क्योंकि उसे $\Zero$ और $\Succ$ से क्रमिक संयोजन
द्वारा परिभाषित किया जा सकता है। यह तब उपयोगी है जब हम आदिम पुनरावर्ती
परिभाषाओं में अचरों का उपयोग करना चाहते हैं। मसलन, यदि हम फलन
$f(x) = 2 \cdot x$ परिभाषित करना चाहें, तो उसे $\fn{const}_n(x)$ और
गुणा से संयोजन द्वारा $f(x) =
\Mult(\fn{const}_2(x), \Proj{1}{0}(x))$ के रूप में पा सकते हैं। आगे से
हम इस युक्ति का उपयोग करेंगे।

\begin{prop}
  घातांक फलन $\fn{exp}(x, y) = x^y$ आदिम पुनरावर्ती है।
\end{prop}

\begin{proof}
  हम $\fn{exp}$ को आदिम पुनरावर्ती रूप से इस प्रकार परिभाषित कर सकते हैं:
  \begin{align*}
    \fn{exp}(x, 0) & = 1\\
    \fn{exp}(x, y+1) & = \Mult(x, \fn{exp}(x,y)).
    \intertext{ठीक-ठीक कहें तो यह आदिम पुनरावर्ती फलनों से पुनरावर्ती
      परिभाषा नहीं है। किंतु औपचारिक रूप से हमारे पास है:}
    \fn{exp}(x, 0) & = f(x)\\
    \fn{exp}(x, y+1) & = g(x, y, \fn{exp}(x,y)).
    \intertext{जहाँ}
    f(x) & = \Succ(\Zero(x)) = 1\\
    g(x, y, z) & = \Mult(\Proj{3}{0}(x, y, z), \Proj{3}{2}(x, y, z)) = x \cdot z
  \end{align*}
  और इस प्रकार $f$ तथा $g$ आदिम पुनरावर्ती फलनों से संयोजन द्वारा
  परिभाषित हैं।
\end{proof}

\begin{prop}
  निम्न प्रकार परिभाषित पूर्ववर्ती फलन $\fn{pred}(y)$
  \[
  \fn{pred}(y) = \begin{cases}
    0 & \text{यदि $y=0$}\\
    y-1 & \text{अन्यथा}
  \end{cases}
  \]
  आदिम पुनरावर्ती है।
\end{prop}

\begin{proof}
 ध्यान दें कि
 \begin{align*}
   \fn{pred}(0) & = 0 \text{ और}\\
   \fn{pred}(y+1) & = y.
 \end{align*}
 यह लगभग आदिम पुनरावर्ती परिभाषा है। ठीक-ठीक कहें तो यह आदिम पुनरावर्तन
 द्वारा परिभाषा के प्रतिरूप में नहीं आती, क्योंकि उस प्रतिरूप को कम-से-कम
 एक अतिरिक्त तर्क~$x$ चाहिए। यह इस अर्थ में भी असामान्य है कि
 $\fn{pred}(y)$ का वास्तव में उपयोग नहीं होता जब $\fn{pred}(y+1)$ को
 परिभाषित किया जाता है। किंतु पहले हम $\fn{pred}'(x, y)$ को इस प्रकार परिभाषित कर सकते हैं:
 \begin{align*}
   \fn{pred}'(x, 0) & = \Zero(x) = 0,\\
   \fn{pred}'(x, y+1) & = \Proj{3}{1}(x, y, \fn{pred'}(x, y)) = y.
 \end{align*}
और फिर उससे संयोजन द्वारा $\fn{pred}$ परिभाषित कर सकते हैं, मसलन
$\fn{pred}(x) = \fn{pred}'(\Zero(x), \Proj{1}{0}(x))$ के रूप में।
\end{proof}

\begin{prop}
  क्रमगुणित फलन $\fn{fac}(x) = \fact{x} = 1 \cdot 2 \cdot 3
  \cdot \dots \cdot x$ आदिम पुनरावर्ती है।
\end{prop}

\begin{proof}
  स्पष्ट आदिम पुनरावर्ती परिभाषा यह है:
  \begin{align*}
    \fn{fac}(0) &= 1\\
    \fn{fac}(y+1) & = \fn{fac}(y) \cdot (y+1).
    \intertext{औपचारिक रूप से हमें पहले एक द्वि-स्थानी फलन $h$ परिभाषित करना है:}
    h(x, 0) & = \fn{const}_1(x)\\
    h(x, y+1) & = g(x, y, h(x, y))
    \intertext{जहाँ $g(x, y, z) = \Mult(\Proj{3}{2}(x, y, z),
      \Succ(\Proj{3}{1}(x, y, z)))$ है, और फिर रखें:}
    \fn{fac}(y) & = h(\Proj{1}{0}(y), \Proj{1}{0}(y)) = h(y,y).
  \end{align*}
  आगे से हम कुछ अधिक अनौपचारिक रहेंगे और संयोजन तथा आदिम पुनरावर्तन
  द्वारा पूर्ण औपचारिक परिभाषाएँ नहीं देंगे।
\end{proof}

\begin{prop}
  निम्न प्रकार परिभाषित छँटा हुआ घटाव $x \tsub y$
  \[
  x \tsub y = \begin{cases}
    0 & \text{यदि $x < y$}\\
    x-y & \text{अन्यथा}
  \end{cases}
  \]
  आदिम पुनरावर्ती है।
\end{prop}

\begin{proof}
  हमारे पास है:
  \begin{align*}
    x \tsub 0 & = x\\
    x \tsub (y+1) & = \fn{pred}(x \tsub y) 
  \end{align*}
\end{proof}

\begin{prop}
 $x$ और $y$ के बीच की दूरी $\left|x-y\right|$ आदिम पुनरावर्ती है।
\end{prop}

\begin{proof}
  हमारे पास $\left| x-y \right| = (x \tsub y) + (y \tsub x)$ है;
  अतः दूरी को $+$ और $\tsub$ से संयोजन द्वारा परिभाषित किया जा सकता है,
  जो आदिम पुनरावर्ती हैं।
\end{proof}

\begin{prop}
  $x$ और $y$ का अधिकतम $\fn{max}(x,y)$ आदिम पुनरावर्ती है।
\end{prop}

\begin{proof}
  हम $\fn{max}(x,y)$ को $+$ और $\tsub$ से संयोजन द्वारा इस प्रकार
  परिभाषित कर सकते हैं:
  \[
  \fn{max}(x,y) \defis x + (y \tsub x).
  \]
  यदि $x$ अधिकतम है, अर्थात् $x \ge y$, तो $y \tsub x = 0$, इसलिए $x
  + (y \tsub x) = x + 0 = x$। यदि $y$ अधिकतम है, तो $y \tsub x =
  y - x$, और इसलिए $x + (y \tsub x) = x + (y - x) = y$।
\end{proof}

\begin{prop}
  \ollabel{prop:min-pr}
  $x$ और $y$ का न्यूनतम $\fn{min}(x,y)$ आदिम पुनरावर्ती है।
\end{prop}

\begin{proof}
  अभ्यास।
\end{proof}

\begin{prob}
  \olref[cmp][rec][exa]{prop:min-pr} सिद्ध करें।
\end{prob}

\begin{prob}
दिखाएँ कि \[
f(x, y) =
2^{(2^{\iddots^{2^{x}}})}\raisebox{1ex}{\bigg\rbrace}
\raisebox{1ex}{\text {$y$ बार $2$}}\] आदिम पुनरावर्ती है।
\end{prob}

\begin{prob}
दिखाएँ कि पूर्णांक भाग $d(x, y) = \lfloor x/y \rfloor$ (अर्थात् ऐसा
भाग जिसमें दशमलव बिंदु के बाद की हर चीज़ छोड़ दी जाती है) आदिम
पुनरावर्ती है। जब $y = 0$ हो, तब हम $d(x, y) = 0$ निर्धारित करते हैं।
आदिम पुनरावर्तन और संयोजन का उपयोग करके~$d$ की स्पष्ट परिभाषा दें।
\end{prob}

\begin{prop}
आदिम पुनरावर्ती फलनों का समुच्चय निम्न दो संक्रियाओं के अधीन संवृत है:
\begin{enumerate}
\item परिमित योग: यदि $f(\vec x, z)$ आदिम पुनरावर्ती है, तो निम्न फलन
भी ऐसा ही है:
\[
g(\vec x, y) \defis \sum_{z = 0}^y f(\vec x, z).
\]
\item परिमित गुणनफल: यदि $f(\vec x, z)$ आदिम पुनरावर्ती है, तो निम्न
फलन भी ऐसा ही है:
\[
h(\vec x, y) \defis \prod_{z = 0}^y f(\vec x, z).
\]
\end{enumerate}
\end{prop}

\begin{proof}
मसलन, परिमित योग निम्न समीकरणों से पुनरावर्ती रूप में परिभाषित होते हैं:
\begin{align*}
  g(\vec x, 0) & = f(\vec x, 0)\\
  g(\vec x, y+1) & = g(\vec x, y) + f(\vec x, y+1).
\end{align*}
\end{proof}


\end{document}
